% problem-000209 2 钪氧化物晶胞计算
\section*{2. 钪氧化物晶胞计算}
过氧化钙 $\left( \mathrm { C a O _ { 2 } } \right)$ 微溶于⽔，溶于酸，可作为⽤医⽤防腐剂和消毒剂，也可作为改良剂为农业、园艺和⽣物技术应⽤提供氧⽓。

\subsection*{2-1}
以下是⼀种实验室制备过氧化钙⽅法的流程图（图2.1）：

（图：）

\subsection*{2-1}
-1 选择步骤①中应当略过量的试剂：（ ） A. 稀盐酸； B. 碳酸钙。

\subsection*{2-1}
-2 解释步骤①中该试剂应当略过量的原因。

\subsection*{2-1}
-3 步骤①中煮沸操作的⽬的是什么？

\subsection*{2-1}
-4 步骤①中过滤操作的⽬的是什么？

\subsection*{2-2}
-1 写出步骤②中⽣成 $\mathrm { C a O } _ { 2 } { \cdot } 8 \mathrm { H } _ { 2 } \mathrm { O }$ 的反应⽅程式

\subsection*{2-2}
-2 步骤②中的混和操作，应当选择哪种⽅式为宜？（ ）A. 将氯化钙溶液滴⼊氨⽔-双氧⽔混合液中；B. 将氨⽔-双氧⽔混合液滴⼊氯化钙溶液中。

\subsection*{2-2}
-3 解释选择2-2-2中混和⽅式的原因。

\subsection*{2-3}
可采⽤量⽓法测定过氧化钙的含量。量⽓法实验装置如图 2.2 所⽰。加热⾄ $3 5 0 ~ ^ { \circ } \mathrm { C }$ 以上时 $\mathrm { C a O }$ ，迅速分解⽣成CaO和 $\mathrm { O } _ { 2 \mathrm { o } }$ 。采⽤量⽓法可以测定产品中 $\mathrm { C a O _ { 2 } }$ 的纯度（假设杂质不产⽣⽓体）。

\subsection*{2-3}
-1 使⽤量⽓管的注意事项中，下列说法正确的是：（ ）A. 使⽤前需要检漏； B. 初始读数时量⽓管与⽔准管液⾯相平；C. 读数时，视线与最低凹液⾯处平⻬； D. 读终体积时，停⽌加热后需⽴即读数；E. 读终体积时，⽓体需先冷却⾄室温； F. 读终体积时，量⽓管与⽔准管液⾯⽆需相平；

（图：）
图 2.2 量⽓法装置图
